\documentclass{article}
\usepackage[utf8]{inputenc}
\usepackage[russian]{babel} % For Russian language support
\usepackage{amsmath,amsfonts,amssymb} % For mathematical symbols

\title{Пример научной статьи}
\author{Автор}
\date{\today}

\begin{document}

\maketitle

\begin{abstract}
Данная статья представляет собой пример документа \LaTeX{} типа article с введением и заключением. В работе рассматриваются основные структурные элементы научной статьи и их оформление в системе \LaTeX{}.
\end{abstract}

\section{Введение}
Научная статья является одним из основных способов представления результатов исследований. В этой статье мы рассмотрим структуру типичной научной статьи, оформленной с помощью системы \LaTeX{}. \LaTeX{} — это мощная система подготовки документов, особенно популярная в научной среде благодаря возможности красивого оформления формул и автоматического создания библиографии.

Введение обычно содержит описание проблемы, цели исследования, краткий обзор существующих работ и структуры статьи. Оно должно заинтересовать читателя и дать ему представление о том, что он найдет в остальной части статьи.

\section{Основная часть}
Основная часть статьи содержит описание методов, результатов и анализ. В этой части обычно приводятся теоретические выкладки, экспериментальные данные и их интерпретация. Структура основной части может варьироваться в зависимости от тематики статьи и требований конкретного журнала или конференции.

Для иллюстрации возможностей \LaTeX{} приведем пример математической формулы:
\begin{equation}
    e^{i\pi} + 1 = 0
\end{equation}

Это тождество Эйлера, связывающее пять фундаментальных математических констант.

\section{Заключение}
В заключении подводятся итоги исследования, формулируются основные выводы и указываются направления для дальнейшей работы. В данной статье мы рассмотрели основные структурные элементы научной статьи и показали, как их можно оформить с помощью системы \LaTeX{}.

Использование \LaTeX{} позволяет создавать профессионально выглядящие документы с высоким качеством типографики. Это особенно важно для научных публикаций, где точность и ясность представления информации имеют первостепенное значение.

% Bibliography can be added here if needed
% \begin{thebibliography}{9}
% \bibitem{latexcompanion} 
% Michel Goossens, Frank Mittelbach, and Alexander Samarin. 
% \textit{The \LaTeX\ Companion}. 
% Addison-Wesley, Reading, Massachusetts, 1993.
% \end{thebibliography}

\end{document}